\subsection{ALGEBRAS}

Let \(\mathbf{M}\) be a unitary module over \(\mathbf{K}\) with a bilinear mapping \((x,y)\mapsto xy\) of \(\mathbf{M}\times \mathbf{M}\) into \(\mathbf{M}\) . All the axioms for algebras are satisfied except associativity of multiplication. By an abuse of language, \(\mathbf{M}\) is called a not necessarily associative algebra over \(\mathbf{K}\) , or sometimes, when no confusion can arise, an algebra over \(\mathbf{K}\) . In this no. we shall use the latter notation.

If the K-module M is given the multiplication \((x,y)\mapsto yx\) an algebra is obtained called the opposite of the above algebra.

A sub-K-module N of M which is stable under multiplication is given the structure of an algebra over K in an obvious way. N is called a subalgebra of M. N is called a left (resp. right) ideal of M if the conditions \(x \in N\) , \(y \in M\) imply \(yx \in N\) (resp. \(xy \in N\) ). If N is both a left ideal and a right ideal of M, N is called a two-sided ideal of M. In this case the multiplication on M enables us to define, on passing to the quotient, a bilinear multiplication on the quotient module M/N such that M/N has an algebra structure. M/N is called the quotient algebra of M by N.

Let \(\mathbf{M}_1\) and \(\mathbf{M}_2\) be two algebras over \(\mathbf{K}\) and \(\phi\) a mapping of \(\mathbf{M}_1\) into \(\mathbf{M}_2\) . \(\phi\) is called a homomorphism if \(\phi\) is K-linear and \(\phi(xy) = \phi(x)\phi(y)\) for \(x\in \mathbf{M}_1\) , \(y\in \mathbf{M}_1\) . The kernel \(\mathbf{N}\) of \(\phi\) is a two-sided ideal of \(\mathbf{M}_1\) and the image of \(\phi\) is a subalgebra of \(\mathbf{M}_2\) . On passing to the quotient, \(\phi\) defines an isomorphism of the algebra \(\mathbf{M}_1 / N\) onto the algebra \(\phi (\mathbf{M}_1)\) .

Let \(\mathbf{M}\) be an algebra over \(\mathbf{K}\) . A mapping \(\mathbf{D}\) of \(\mathbf{M}\) into \(\mathbf{M}\) is called a derivation of \(\mathbf{M}\) if it is \(\mathbf{K}\) -linear and \(\mathrm{D}(xy) = (\mathrm{D}x)y + x(\mathrm{D}y)\) for all \(x \in \mathbf{M}\) and \(y \in \mathbf{M}\) . This definition generalizes Definition 3 of Algebra, Chapter IV, § 4, no. 3. The kernel of a derivation of \(\mathbf{M}\) is a subalgebra of \(\mathbf{M}\) . If \(\mathrm{D}_1\) and \(\mathrm{D}_2\) are derivations of \(\mathbf{M}\) , then \(\mathrm{D}_1\mathrm{D}_2 - \mathrm{D}_2\mathrm{D}_1\) is a derivation of \(\mathbf{M}\) (cf.~Algebra, Chapter IV, § 4, no. 3, Proposition 5: the proof of this proposition does not use the associativity of the algebra).

Let \(\mathbf{M}_1\) and \(\mathbf{M}_2\) be two algebras over K. On the product K-module \(\mathbf{M} = \mathbf{M}_1 \times \mathbf{M}_2\) we define a multiplication by writing

\[
(x _ {1}, x _ {2}) (y _ {1}, y _ {2}) = (x _ {1} y _ {1}, x _ {2} y _ {2}),
\]

for all \(x_1, y_1\) in \(\mathbf{M}_1, x_2, y_2\) in \(\mathbf{M}_2\) . The algebra thus defined is called the product algebra of \(\mathbf{M}_1\) and \(\mathbf{M}_2\) . The mapping \(x_1 \mapsto (x_1, 0)\) (resp. \(x_2 \mapsto (0, x_2)\) ) is an isomorphism of \(\mathbf{M}_1\) (resp. \(\mathbf{M}_2\) ) onto a two-sided ideal of M. Under these isomorphisms \(\mathbf{M}_1\) and \(\mathbf{M}_2\) are identified with two-sided ideals of M. The K-module M is then the direct sum of \(\mathbf{M}_1\) and \(\mathbf{M}_2\) . Conversely, let M be an algebra over K and \(\mathbf{M}_1, \mathbf{M}_2\) two two-sided ideals of M such that M is the direct sum of \(\mathbf{M}_1\) and \(\mathbf{M}_2\) . Then \(\mathbf{M}_1\mathbf{M}_2 \subset \mathbf{M}_1 \cap \mathbf{M}_2 = \{0\}\) ; then, if \(x_1, y_1\) belong to \(\mathbf{M}_1\) and \(x_2, y_2\) to \(\mathbf{M}_2\) , then \((x_1 + x_2)(y_1 + y_2) = x_1y_1 + x_2y_2\) , so that M is identified with the product algebra \(\mathbf{M}_1 \times \mathbf{M}_2\) . Every left (resp. right, two-sided) ideal of \(\mathbf{M}_1\) is a left (resp. right, two-sided) ideal of M. We leave to the reader the task of formulating the analogous results in the case of an arbitrary finite family of algebras.

Let \(\mathbf{M}\) be an algebra over \(\mathbf{K}\) and suppose that the \(\mathbf{K}\) -module \(\mathbf{M}\) admits a basis \((a_{\lambda})_{\lambda \in \mathbf{L}}\) . There exists a unique system \((\gamma_{\lambda \mu \nu})_{(\lambda, \mu, \nu) \in \mathbf{L} \times \mathbf{L} \times \mathbf{L}}\) of elements of \(\mathbf{K}\) such that \(a_{\lambda}a_{\mu} = \sum_{\nu} \gamma_{\lambda \mu \nu}a_{\nu}\) for all \(\lambda, \mu\) in \(\mathbf{L}\) . The \(\gamma_{\lambda \mu \nu}\) are called the constants of structure of \(\mathbf{M}\) with respect to the basis \((a_{\lambda})\) .

Let \(\mathbf{M}\) be an algebra over \(\mathbf{K}\) , \(\mathbf{K}_0\) a commutative ring with unit element and \(\rho\) a homomorphism of \(\mathbf{K}_0\) into \(\mathbf{K}\) mapping unit element to unit element. Then \(\mathbf{M}\) can be considered as an algebra over \(\mathbf{K}_0\) by writing \(\alpha.x = \rho(\alpha).x\) for \(\alpha \in \mathbf{K}_0, x \in \mathbf{M}\) . This is the case in particular when \(\mathbf{K}_0\) is a subring of \(\mathbf{K}\) containing the unit element and \(\rho\) is taken to be the inclusion mapping of \(\mathbf{K}_0\) into \(\mathbf{K}\) .

Let M be an algebra over K, \(K_{1}\) a commutative ring with unit element and \(\sigma\) a homomorphism of K into \(K_{1}\) mapping unit element to unit element. Let \(M_{(K_{1},\sigma)} = M_{(K_{1})}\) be the \(K_{1}\) -module derived from M by extending the ring of scalars to \(\mathbf{K}_1\) (Algebra, Chapter II, § 5). The product on M defines canonically a \(\mathbf{K}_1\) -bilinear mapping of \(\mathrm{M}_{(\mathbf{K}_1)} \times \mathrm{M}_{(\mathbf{K}_1)}\) into \(\mathrm{M}_{(\mathbf{K}_1)}\) (Algebra, Chapter IX, § 1, no. 4) such that \(\mathrm{M}_{(\mathbf{K}_1)}\) is given the structure of an algebra over \(\mathbf{K}_1\) (which is said to be derived from M by extending the ring of scalars to \(\mathbf{K}_1\) ). This is the case in particular when K is a subring of \(\mathbf{K}_1\) containing the unit element and \(\sigma\) the inclusion mapping of K into \(\mathbf{K}_1\) .

